\section{Three further quartic reflection-field class certificates}
\label{sec:cubic-three-quartic-class-certificates}

Put $p=269$, $t=\sqrt[3]{p}\in\mathbb R$, $K=\mathbb Q(\sqrt{-3})$,
and $z=\zeta_5$.  For $d\in\{-11,2,-22\}$ define
\[
 w_d=\sqrt d\,(z-z^{-1}),\qquad F_d=\mathbb Q(w_d),\qquad
 L_d=\mathbb Q(t,w_d),\qquad M_d=KL_d.
 \tag{TQC.1}
\]
Let $D_d$ be the fundamental discriminant of $\mathbb Q(\sqrt d)$,
so $D_{-11}=-11$, $D_2=8$, and $D_{-22}=-88$.
Let $\psi$ be the quartic character modulo $5$ with $\psi(2)=i$,
put $\delta=\psi^2$, and put $\chi_d=\psi\tau_{D_d}$.
All Dirichlet characters are extended by zero at nonunits.

\begin{theorem}[AC-G135-quartic-classes: three exact class sectors]
\label{AC-G135-quartic-classes}
For each of the three values of $d$ in (TQC.1),
\[
       \operatorname{Cl}(L_d)[269]=0.
 \tag{TQC.2}
\]
Consequently the nontrivial cubic-character sector of
$\operatorname{Cl}(M_d)[269]$ is zero.  This conclusion does not
assert a complete class number, a complete unit basis, vanishing
of the trivial cubic sector of $\operatorname{Cl}(M_d)[269]$,
or any subsequent deformation-host or model construction.
\end{theorem}

\begin{proof}
\proofstep{1. Exact fields, discriminants, signatures, and roots of unity.}
The fields $\mathbb Q(z)$ and $\mathbb Q(\sqrt d)$ are disjoint,
since the former has unique quadratic subfield $\mathbb Q(\sqrt5)$.
In their compositum let $s$ send $z$ to $z^2$ and fix $\sqrt d$,
and let $r$ negate $\sqrt d$ and fix $z$.
The stabilizer of $w_d$ is $\{1,s^2r\}$, exactly the kernel of
$\psi\tau_{D_d}$ on $C_4\times C_2$.
Thus $F_d/\mathbb Q$ is cyclic of degree four, with character group
$\{1,\delta,\chi_d,\bar\chi_d\}$, and
\[
 w_d^4+5d w_d^2+5d^2=0,\qquad
 w_d^2=-d(5\pm\sqrt5)/2.
 \tag{TQC.3}
\]
The character $\chi_d$ is primitive of conductor $5|D_d|$.
It is even for $d<0$ and odd for $d=2$.
The conductor--discriminant formula for this abelian field gives
$D(F_d)=5^3D_d^2$.  In particular $F_d$ is unramified at $3$ and $p$.

The field $T=\mathbb Q(t)$ has discriminant $-3p^2$.
Here is the local check of the only possible index issue.
At $p$ the polynomial $X^3-p$ is Eisenstein and tamely ramified.
At $3$, the equation $v+3v^2+3v^3=-30$ has a solution in
$\mathbb Z_3$, by its unit derivative; hence $-(1+3v)$ is a
cube root of $p$ in $\mathbb Q_3$.
The remaining quadratic factor generates $\mathbb Q_3(\sqrt{-3})$.
The local discriminant exponents are therefore $2$ at $p$ and $1$
at $3$, with no other primes by the polynomial discriminant.
The degrees three and four imply $T\cap F_d=\mathbb Q$.
Their discriminants are coprime.  Locally at each rational prime,
one integral algebra is unramified, and its tensor product with
the other maximal order is finite etale over a product of DVRs,
hence normal.  Thus the tensor order is maximal, and
\[
 D(L_d)=(3p^2)^4 5^9|D_d|^6,
 \qquad
 (r_1,r_2,r)=
 \begin{cases}(4,4,7),&d=-11,-22,\\(0,6,5),&d=2.
 \end{cases}
 \tag{TQC.4}
\]
The signatures follow either from (TQC.3), or from the one real
and one conjugate pair of embeddings of $T$.

The splitting field of $X^3-p$ has group $S_3$ and unique quadratic
subfield $K$.  Its intersection with the abelian field $F_d$ is
trivial, since $K$ is ramified at $3$ and $F_d$ is not.
Consequently $\operatorname{Gal}(M_d/\mathbb Q)=S_3\times C_4$,
and $L_d$ is fixed by one reflection in $S_3$.
That reflection normally generates the entire $S_3$ factor, so
every abelian subextension of $L_d$ lies in $F_d$.
All roots of unity in $L_d$ therefore lie in $F_d$.
For $d<0$ this quartic field is real.  For $d=2$ it has the unique
quadratic subfield $\mathbb Q(\sqrt5)$; the cyclotomic degrees
dividing four exclude every root of unity other than $\pm1$,
except possibly fifth roots, which would give
$F_2=\mathbb Q(\zeta_5)$, contrary to its ramification at $2$.
Thus in all three cases
\[
                         |\mu_{L_d}|=2.
 \tag{TQC.5}
\]

\proofstep{2. Exact global units and a prime-to-269 index.}
All finite input data are rational.  The ordered basis is
$t^a w_d^b$, $0\le a<3$, $0\le b<4$, with array index $4a+b$.
For each field, the file \texttt{all\_units\_exact\_certificate.json}
gives respectively seven, five, or seven rational arrays, their
common denominators, and monic integral annihilating polynomials
of degree twelve with constant coefficient $\pm1$.
The three exact-file hashes, in the order $d=-11,2,-22$, are
\begin{gather*}
\texttt{59191fdbe9e6b85e9efa69fbf51968bce3}\allowbreak
\texttt{d728994d89da5ba9768cf14ba495b1},\\
\texttt{ef0714465df3215cc3607be89c37deab27}\allowbreak
\texttt{1c6fe74ee1babe50c07e6225b2adb1},\\
\texttt{d1f659cb8ad5792f5d0fe48196d9206934}\allowbreak
\texttt{e54afd333d1d9a06c2615b7766c428}.
\end{gather*}
The arrays and polynomials are part of the retained finite
certificate, not numerical approximations or an appeal to a
provisional unit-basis assertion.

The exact verifier multiplies two arrays in $\mathbb Q[t,w_d]$,
reducing only by $t^3=p$ and (TQC.3), and evaluates each displayed
polynomial by Horner's rule.  Every value is exactly zero.
A monic integral relation with constant coefficient $\pm1$
proves both the element and its inverse integral.
Thus all nineteen elements are actual global units.
The candidate coordinate was $w_d-t$; its displayed inverse maps
to $t,w_d$ are checked in the same rational algebra.
A separate raw-polynomial verifier checks these maps, all
annihilations, and the conversion to the tensor arrays.
No complete unit-basis or class-group claim from the candidate
extraction is used.

For saturation take $\ell=40351$, primitive root $3$, and
$h=3^{150}=10281\pmod\ell$, of order $269$.
Trial division proves $\ell$ prime; the powers for $2,3,5,269$
prove that $3$ is primitive.
There are three roots of $t^3=p$ and four roots of (TQC.3)
modulo $\ell$.  All twelve pairs have nonzero derivatives.
They lift to embeddings $L_d\hookrightarrow\mathbb Q_\ell$;
algebraic integers map to $\mathbb Z_\ell$, so these are actual
residue maps of the full integer ring.
Each unit denominator is invertible modulo $\ell$.
For its image $\bar u_i$ under the $j$th map, write
$\bar u_i^{150}=h^{c_{ji}}$.
The complete twelve-row matrices occur in
\texttt{269\_saturation\_certificate.json}.
Their first $r$ rows have exact integer determinants
\[
\begin{array}{c|r|r}
d&\det(c_{ji})_{1\le i,j\le r}&\det\bmod269\\\hline
-11&-15656332564584192&194\\
2&30877898400&106\\
-22&102858999958416&115.
\end{array}
\tag{TQC.6}
\]
Exact modular powers verify every matrix entry.
An independent implementation recomputes all thirty-six residue
rows and the determinants by the Leibniz formula.

The resulting homomorphism from global units to
$\mathbb F_{269}^{12}$ kills $269$th powers and all torsion.
The matrix has rank $r$.  By the unit theorem,
$\mathcal O_{L_d}^{\times}/\mu_{L_d}$ is free of rank $r$.
The subgroup $U_d$ generated by the displayed units and $-1$
therefore has full rank and index
\[
 m_d=[\mathcal O_{L_d}^{\times}:U_d],\qquad 269\nmid m_d.
 \tag{TQC.7}
\]
Indeed full rank follows already from independence modulo $269$;
a nonzero rational relation could be divided until its integer
coefficients were not all divisible by $269$, contradicting
(TQC.6).  The same matrix says that the induced map of the two
rank-$r$ lattices is invertible modulo $269$, proving the assertion
about the index.

\proofstep{3. The exact projected regulators.}
Put $\rho=p^{1/3}$, $\omega=e^{2\pi i/3}$, and
$a_d=\sqrt{|d|(5+\sqrt5)/2}$,
$b_d=\sqrt{|d|(5-\sqrt5)/2}$.
For $d<0$, take the four real embeddings
$(t,w)=(\rho,\pm a_d),(\rho,\pm b_d)$, with weight one,
and the four complex-place representatives
$(\rho\omega,\pm a_d),(\rho\omega,\pm b_d)$, with weight two.
For $d=2$, take the six complex representatives
\[
 (t,w)=(\rho\omega^j,ia_d),(\rho\omega^j,ib_d),
                          \qquad j=0,1,2,
 \tag{TQC.8}
\]
all with weight two.
If $q_v$ is that weight, form the matrix
$q_v\log|\sigma_v(u_i)|$, omit its last row, and take the
absolute determinant $R_{U_d}$.
This is the projected weighted-log regulator, so
\[
                         R_{U_d}=m_d R_{L_d}.
 \tag{TQC.9}
\]
No Euclidean hyperplane-covolume factor is inserted.
All rows are kept in the certificate, every column sum encloses
zero, and the ball determinant is strictly positive.
Evaluation uses respectively $4000,6000,12000$ bits, before the
lower-precision analytic calculation.  It yields
\[
\begin{array}{c|c}
d&\text{strict enclosing interval for }R_{U_d}\\\hline
-11&(4789500350,4789500351)\\
2&(3666775011,3666775012)\\
-22&(6035733830500,6035733830501).
\end{array}
\tag{TQC.10}
\]
These intervals are calculated from the exact arrays, not used
as input bounds; in particular the large cancellation in the
last field is included in the outward-rounded determinant.

\proofstep{4. A denominator bound from ideal counting.}
We give the isolation argument without a regulator lower bound.
For any degree-$n$ field of absolute discriminant $D$ and $r_2$
complex pairs, define exact integers
\[
 U=\left\lceil\frac{n!4^{r_2}(\lfloor\sqrt D\rfloor+1)}
                         {n^n3^{r_2}}\right\rceil,
 \qquad H=4^n\left\lceil\sqrt[3]{U^4}\right\rceil.
 \tag{TQC.11}
\]
The Minkowski ideal-class bound, with $\pi>3$, supplies an
integral representative of every class of norm at most $U$.
For $s=4/3$, elementary ideal counting and Euler factors give
\[
 h\le\#\{\mathfrak a:N\mathfrak a\le U\}
 \le U^{4/3}\zeta_L(4/3)
 \le U^{4/3}\zeta(4/3)^n
 \le4^n U^{4/3}\le H.
 \tag{TQC.12}
\]
There are at most $n$ primes over each rational prime and each
residue degree is at least one, proving the middle inequality.
The final analytic estimate is simply
$\sum_{m\ge2}m^{-4/3}\le\int_1^\infty x^{-4/3}dx=3$.
For (TQC.4), the exact $(U,H)$ are
\[
\begin{array}{c|r|r}
d&U&H\\\hline
-11&14883725864432&6142362534283955232833536\\
2&10178435786078&3700791495413757380657152\\
-22&7620467642588970&25159116940426138598695239680.
\end{array}
\tag{TQC.13}
\]
If an interval encloses $r=h/m>1$ and proves
$|r-c|<1/H$ for an integer $c$, then $m<h\le H$.
The nonzero integer $h-cm$ would force
$|r-c|\ge1/m>1/H$.  Thus that strict interval proves $h=cm$.
This uses the actual rational value $h/m$, not a rounded estimate.

\proofstep{5. Artin factorization and exact Dirichlet constants.}
Let $\varrho$ be the standard two-dimensional representation of
the $S_3$ splitting field of $X^3-p$.
The permutation representation for $L_d$ is
$(1\oplus\varrho)\otimes(1\oplus\delta\oplus\chi_d\oplus\bar\chi_d)$.
Artin formalism, with its ramified inertia-invariant Euler factors,
therefore gives
\[
 \operatorname*{res}_{s=1}\zeta_{L_d}(s)
 =L(1,\delta)|L(1,\chi_d)|^2
       L(1,\varrho)L(1,\varrho\delta)|L(1,\varrho\chi_d)|^2.
 \tag{TQC.14}
\]
The quadratic character here is the primitive conductor-five
character $\delta$.  It is not the imprimitive square of the
zero-extended function modulo $5|D_d|$: no Euler factor at a prime
of $D_d$ is deleted from $L(1,\varrho\delta)$.

Put $\varphi=(1+\sqrt5)/2$.
The finite Fourier identity modulo five gives
\[
 L(1,\delta)=-\frac1{\sqrt5}\sum_{a=1}^{4}\delta(a)
             \log(2\sin(\pi a/5))=\frac{2\log\varphi}{\sqrt5}.
 \tag{TQC.15}
\]
It follows by multiplying
$-\log|1-re^{i\theta}|=\sum_{n\ge1}r^n\cos(n\theta)/n$
by $\delta(a)$ and summing, then taking $r\uparrow1$ by Abel's
theorem and Dirichlet convergence.  The finite Gauss constant is
$2\cos(2\pi/5)-2\cos(4\pi/5)=\sqrt5>0$.
The last equality uses $\sin(2\pi/5)/\sin(\pi/5)=\varphi$.

For $d=2$, the primitive odd character has conductor $q=40$ and
the exact integer sum gives
\[
 B_{1,\chi_2}=\frac1{40}\sum_{a=1}^{39}a\chi_2(a)=-2i,
 \qquad |L(1,\chi_2)|^2=\frac{\pi^2}{10}.
 \tag{TQC.16}
\]
The second equality is the primitive odd Dirichlet functional
equation and the Gauss norm, or equivalently its finite Fourier
special-value formula.
For $d=-11,-22$, the character is primitive even, of conductor
$q=55,440$ respectively.  Set
\[
 S_d=\sum_{a=1}^{q-1}\chi_d(a)\log(2\sin(\pi a/q)).
 \tag{TQC.17}
\]
The same Abel-limit calculation gives
$S_d=-\mathfrak g(\chi_d)L(1,\bar\chi_d)$:
the cosine sum equals the exponential sum since $\chi_d(-1)=1$.
The primitive Gauss norm $|\mathfrak g(\chi_d)|^2=q$ consequently gives
\[
                         |L(1,\chi_d)|^2=|S_d|^2/q.
 \tag{TQC.18}
\]
No Gauss-sum phase or logarithm branch is guessed.
Every sine in (TQC.17) is strictly positive; the interval calculation
uses its real logarithm and the exact character values $0,\pm1,\pm i$.

\proofstep{6. Actual Hecke conductors and coefficient sequences.}
Write $L(s,\varrho)=\sum a_n n^{-s}$.
Its coefficients are multiplicative.  For a prime $\ell\ne3,p$,
$a_\ell=\#\{x\bmod\ell:x^3=p\}-1$ and
\[
 a_{\ell^{v+1}}=a_\ell a_{\ell^v}
                   -(-3/\ell)a_{\ell^{v-1}},\qquad
 a_1=1,\quad a_{3^v}=1,\quad a_{p^v}=0\quad(v\ge1).
 \tag{TQC.19}
\]
For $\ell\equiv2\pmod3$ the local coefficients are one in even
degree and zero in odd degree.  For $\ell\equiv1\pmod3$, the
test $p^{(\ell-1)/3}\bmod\ell$ decides whether there are three
or zero roots.  The factor at $3$ follows from the split local
cubic algebra above; the factor at $p$ follows from nontrivial
cubic inertia with no standard invariant.
The two twisted sequences are $a_n\delta(n)$ and $a_n\chi_d(n)$.
At every new ramified twist prime the original representation is
unramified, so extension by zero gives the actual primitive factor.

Write $\varrho=\operatorname{Ind}_{K}^{\mathbb Q}\nu$ for its cubic
Kummer character.  This character has conductor $(p)$ over $K$;
the extension is tame at $p$, and $p$ is a cube in $\mathbb Q_3$.
The primes $2,5,11,p$ are inert in $K$.
The norm pullbacks of the primitive characters at $5$ and $11$
retain their conductor exponent one, by the surjective residue
field norm.  At $2$, the only new case is the conductor-eight
quadratic character.  In the unramified quadratic local ring,
\[
 N(1+4a)\equiv1+4\operatorname{Tr}_{\mathbb F_4/\mathbb F_2}(\bar a)
                         \pmod8.
 \tag{TQC.20}
\]
The residue trace is onto, and $\tau_8(5)=-1$.
Thus the pullback is nontrivial on $1+4\mathcal O_{K_2}$ and
trivial on $1+8\mathcal O_{K_2}$, exactly exponent three.
Other local twist factors there are unramified, so cannot cancel it.
At $p$ these three Dirichlet twists are unramified and leave the
cubic conductor unchanged.  At $3$ they are unramified.
The rational conductors are therefore exactly
\[
 Q_0=3p^2=217083,\qquad Q_5=75p^2=5427075,
 \qquad Q_d=75p^2|D_d|^2.
 \tag{TQC.21}
\]
Indeed induction multiplies the norm of the Hecke conductor by
$|D_K|=3$.  The primitive Hecke moduli are
$f=p,5p,5p|D_d|$, respectively.

\proofstep{7. Prime-power finite Fourier transform and Poisson normalization.}
The imaginary quadratic Minkowski bound for $K$ is
$2\sqrt3/\pi<2$, so every ideal class has an integral ideal of
norm one; hence $h(K)=1$.
Each Hecke character above is a primitive finite multiplicative
character $\nu_f$ on $\mathcal O_K/f\mathcal O_K$, trivial on the
six units and extended by zero on nonunits.
For the cubic character this also follows from $9\mid p+1$ and
its odd order; norm pullbacks are trivial on these units.
Put
\[
 a=\frac{2\pi}{\sqrt3 f}=\frac{2\pi}{\sqrt Q},\qquad
 \Theta(t)=\frac16\sum_{\alpha\in\mathcal O_K}
        \nu_f(\alpha)e^{-a|\alpha|^2t}
        =\sum_{n\ge1}b_n e^{-ant}.
 \tag{TQC.22}
\]
The last equality counts the six generators of each ideal.

Here the finite Fourier argument includes prime powers.
The additive pairing
\[
 (\alpha,\beta)\longmapsto
 \exp\left(2\pi i\operatorname{Tr}_{K/\mathbb Q}
               \frac{\alpha\bar\beta}{f\sqrt{-3}}\right)
 \tag{TQC.23}
\]
is perfect on $\mathcal O_K/f\mathcal O_K$ by its inverse-different
description.  In particular it is primitive at $2$.
At a unit frequency, substitution gives Fourier sum
$G\bar\nu_f(\bar\beta)$, where $G$ is the frequency-one Gauss sum.
At a nonunit frequency at a factor $\mathfrak q^a$ of the modulus,
choose a unit $u$ in $1+\mathfrak q^{a-1}$ on which the character
is nontrivial if $a\ge2$; for $a=1$, choose a residue unit on which
it is nontrivial.  Primitivity supplies it, and the Chinese remainder
theorem makes it one at all other factors.
Then $(u-1)\bar\beta=0$ modulo the conductor, so multiplication
by $u$ fixes the additive character but multiplies the sum by
$\nu_f(u)\ne1$.  Thus the sum vanishes.
Finite additive Parseval, with $f^2$ residue classes, now gives
$|G|^2=f^2$.  This proof does not assume the modulus squarefree.

The lattice $f\mathcal O_K\subset\mathbb C$ has area
$f^2\sqrt3/2$ and Euclidean dual
$2/(f\sqrt{-3})\mathcal O_K$.
If $b=2/\sqrt{-3}$, then
\[
 \operatorname{Re}(\alpha\overline{b\beta})
       =\operatorname{Tr}_{K/\mathbb Q}
                        (-\alpha\bar\beta/\sqrt{-3}).
 \tag{TQC.24}
\]
The factor two in the Euclidean dual cancels the factor one-half
in the real-part formula.  There is no assertion that two is a unit
modulo $f$.
The Fourier transform of $e^{-at|x|^2}$ is
$\pi/(at)e^{-\pi^2|\xi|^2/(at)}$.
Applying Poisson to each residue coset, (TQC.23)--(TQC.24) give
the scalar $1/(ft)$, the exponent $a|\beta|^2/t$, and the inverse
character at $\bar\beta$.  Conjugating $\beta$ preserves its norm.
Consequently
\[
 \Theta(t)=\varepsilon t^{-1}\bar\Theta(1/t),\qquad
                   \varepsilon=G/f,\quad |\varepsilon|=1.
 \tag{TQC.25}
\]
Both constant terms are zero.  The Mellin integral converges at
both ends.  For $\operatorname{Re}s>1$, absolute convergence of
$\sum d(n)n^{-s}$ gives its identity with
$\Gamma(s)a^{-s}L(s)$; the exponential bounds at both ends continue
this identity to $s=1$.  Splitting that integral at one and then
substituting $t\mapsto1/t$ in its lower half gives
\[
 L(1)=\sum_{n\ge1}\frac{b_n e^{-an}}n
       +\varepsilon a\sum_{n\ge1}\bar b_n E_1(an),
 \qquad E_1(x)=\int_1^\infty e^{-xt}\frac{dt}{t}.
 \tag{TQC.26}
\]
Every series here is absolutely exponentially convergent.
The computed ball for $\Theta(1)$ excludes zero in every case,
so its actual phase is
\[
                    \varepsilon=\Theta(1)/\overline{\Theta(1)}.
 \tag{TQC.27}
\]
For the two real sequences it is one.  For the complex sequence
the quotient is taken in inclusion arithmetic; no root-number
convention or separately rounded phase is used.

\proofstep{8. Certified tails and class-formula intervals.}
The coefficients in (TQC.19) satisfy $|a_{\ell^v}|\le v+1$,
so every twisted coefficient satisfies $|b_n|\le d(n)\le n$.
For $z_a=e^{-a}$ and $H_N=e^{-a(N+1)}$, the three omitted tails
in the first sum of (TQC.26), its second sum before multiplication
by $a$, and the theta sum are bounded respectively by
\[
 \frac{H_N}{1-z_a},\qquad
 \frac{H_N}{a(1-z_a)},\qquad
 \frac{H_N((N+1)-Nz_a)}{(1-z_a)^2}.
 \tag{TQC.28}
\]
The middle bound uses $E_1(x)\le e^{-x}/x$.
The script accumulates real and imaginary components using exact
integer character values.  Each component receives the full positive
tail radius, which safely encloses the omitted complex series.
It uses $256$ bits for these sums and for the finite Dirichlet sums,
with cutoffs
\[
              N_0=10000,\qquad N_5=40000,
                         \qquad N_d=40000|D_d|.
 \tag{TQC.29}
\]
The high-precision regulator balls from Step 3 are retained when
the analytic working precision is lowered.

The analytic class-number formula, with (TQC.4)--(TQC.5), is
\[
 h_{L_d}R_{L_d}
   =\frac{2\sqrt{D(L_d)}}{2^{r_1}(2\pi)^{r_2}}
                   \operatorname*{res}_{s=1}\zeta_{L_d}(s).
 \tag{TQC.30}
\]
This formula uses $2^4(2\pi)^4$ in the two mixed-signature cases,
and $(2\pi)^6$ in the totally imaginary case.
Combining (TQC.14)--(TQC.30) with the actual determinant gives
the following strict comparisons for $r_d=h_{L_d}/m_d$:
\[
\begin{array}{c|r|c|c}
d&c_d&|r_d-c_d|\text{ is less than}&1/H\text{ is greater than}\\\hline
-11&378&3\cdot10^{-38}&1.62\cdot10^{-25}\\
2&168&6\cdot10^{-38}&2.70\cdot10^{-25}\\
-22&336&3\cdot10^{-37}&3.97\cdot10^{-29}.
\end{array}
\tag{TQC.31}
\]
Each interval also proves $r_d>1$.
The script compares ball objects, not the displayed decimal bounds.
Its exact SHA256 is
\texttt{f6a077327ba0d7dd9a717b82fc6cfd8ee}\allowbreak
\texttt{9c60afacb5c6c7840ac5cc07a230cb0}.
The three complete \texttt{analytic\_interval.json} files retain
all inputs, precision settings, logarithm rows, theta and phase balls,
tail radii, final strict comparisons, timestamps and input hashes.
The integers $378,168,336$ are merely the preselected targets
until this interval and rational-isolation argument is applied.

\proofstep{9. Isolation and reflection descent.}
Step 4 applied to (TQC.31) proves
\[
 h_{L_{-11}}=378m_{-11},\qquad
 h_{L_2}=168m_2,\qquad h_{L_{-22}}=336m_{-22}.
 \tag{TQC.32}
\]
None of the three integers is divisible by $269$, and (TQC.7)
shows that none of the indices is either.  This proves (TQC.2).
It does not prove that an index equals one.

Let $s$ be the reflection fixing $L_d$ in $M_d$.
On ideal classes, norm and extension have composites $2$ and $1+s$.
Since $\operatorname{Cl}(L_d)[269]=0$, the action of $1+s$ on
the full $269$-primary subgroup of $\operatorname{Cl}(M_d)$ is zero.
Thus $s$ acts as $-1$, as do all conjugate reflections.
Their products generate the cubic rotations, which consequently
act as one.  The nontrivial cubic-character sector is therefore zero.
No conclusion about the remaining trivial cubic sector is used.
\end{proof}

\paragraph{Source, computation, and geometric scope.}
The retained analytic class-number formula and projected weighted-log
normalization are A.~Sutherland, MIT 18.785 Lecture 19 (2019), printed
pp.1 and 4, at
\url{https://math.mit.edu/classes/18.785/2019fa/LectureNotes19.pdf}.
The primitive Dirichlet special-value and functional-equation input
is DLMF 25.15.3 and 25.15.5--6, \url{https://dlmf.nist.gov/25.15}.
The Minkowski ideal bound is the same classical input retained in
AC-K19; the new ideal-count isolation proof is given in Step 4.
The prime-power Fourier cancellation, conductor-eight norm calculation,
Poisson constants, phase recovery, absolute tails and reflection
descent are proved explicitly above.
The runtime is python-flint 0.9.0 with its documented outward ball
operations and order-one \texttt{arb.expint}:
\url{https://python-flint.readthedocs.io/en/latest/arb.html}.
No proof of the implementation or new proof of the deep imported
class-number formula is asserted.  No provisional PARI class group,
GRH assumption, complete unit basis, or Lean execution enters the result.

Under the retained class-field and finite-etale-cover dictionary
(SGA1, Expose V, sections 4 and 7), the conclusion controls the
specified cubic sector of unramified abelian covers of
$\operatorname{Spec}\mathcal O_{M_d}$.
Any application to a growing host, changed residual tuple or model
requires its separate arithmetic and local-condition comparison.
